% Generated from locale/tr/content/incompleteness/incompleteness-provability/introduction.tex; do not hand-edit.


\olfileid[tr]{inc}{inp}{int}

\olsection{Giriş}

Hilbert, doğal sayılar gibi bir matematiksel yapı için kurulmuş bir aksiyom
sisteminin, yapı hakkındaki bütün doğru önermelerin türetilmesine izin
vermediği sürece yetersiz olduğunu düşünüyordu. Daha sonra biçimsel türetim
sistemlerine duyduğu ilgiyle birlikte ele alındığında bu görüş, örneğin doğal
sayılar hakkında akıl yürütmek için kullandığımız biçimsel sistemlerin yalnızca
tutarlı değil, aynı zamanda \emph{tam} olmasını da güvence altına almamız
gerektiğini düşündüğünü gösterir: sistemin dilindeki her önerme ya kendisi ya
da değillemesi !!{derivable} olmalıdır. G\"odel'in birinci eksiklik teoremi
böyle bir aksiyom sisteminin bulunmadığını gösterir: aritmetik için tam,
tutarlı ve !!{axiomatizable} hiçbir biçimsel sistem yoktur. Hatta ``yeterince
güçlü'' hiçbir tutarlı, !!{axiomatizable} matematiksel kuram tam değildir.

Hilbert'in daha önemli bir amacı, modern (``klasik'') matematiği temellendirme
programının odağında yer alan, klasik akıl yürütmeyi temsil eden biçimsel
sistemler için sonlu yöntemlere dayanan tutarlılık ispatları bulmaktı. Bu
bakımdan G\"odel'in ikinci eksiklik teoremi Hilbert programına çok daha ağır
bir darbe indirdi. İkinci eksiklik teoremi, birincisi gibi, yaklaşık terimlerle
ifade edilebilir. Kabaca, yeterince güçlü hiçbir aritmetik kuramının kendi
tutarlılığını ispatlayamayacağını söyler. Burada ``yeterince güçlü'' olmayı
$\Th{Q}$'dan biraz daha fazlasını içerecek biçimde anlamamız gerekecek.

G\"odel'in eksiklik teoremine ilişkin özgün ispatının ardındaki fikir
Epimenides paradoksunda bulunabilir. Giritli Epimenides bütün Giritlilerin
yalancı olduğunu ileri sürmüştü; paradoksun daha doğrudan bir biçimi ``bu
tümce yanlıştır'' önermesidir. G\"odel, doğruluğun yerine !!{derivability}
koyarak !!a{sentence} biçimselleştirebildi; bu tümce dolaylı bir yoldan
kendisinin !!{derivable} olmadığını söyler. Bu !!{sentence} !!{derivable} olsaydı
kuram tutarsız olurdu. G\"odel, ele aldığı aksiyom sisteminden bu
!!{sentence}'nin değillemesinin de !!{derivable} olmadığını gösterdi. (Bu
ikinci kısım için G\"odel, $\Th{T}$ kuramının ``$\omega$-tutarlı'' denen
özelliğe sahip olduğunu varsaymak zorundaydı. $\omega$-tutarlılık,
tutarlılıkla ilişkili fakat ondan daha güçlü bir özelliktir.\footnote{Yani her
$\omega$-tutarlı kuram tutarlıdır, fakat bunun tersi doğru değildir.}
G\"odel'den birkaç yıl sonra Rosser, $\Th{T}$'nin yalnızca tutarlı olduğunu
varsaymanın yeterli olduğunu gösterdi.)

İlk güçlük, kendisine gönderme yapan !!a{sentence}'nin nasıl kurulabileceğini
anlamaktır. $\Th{Q}$ dilindeki her $!A$ !!{formula} için $\gn{!A}$,
$\Gn{!A}$ sayısına karşılık gelen rakamı göstersin. Bunun ne anlama geldiğini
düşünün: $!A$, $\Th{Q}$ dilinde !!a{formula}; $\Gn{!A}$ bir doğal sayı;
$\gn{!A}$ ise $\Th{Q}$ dilinde bir \emph{terim}dir. Dolayısıyla $\Th{Q}$
dilindeki her $!A$ !!{formula}, yine $\Th{Q}$ dilinde bir terim olan
$\gn{!A}$ \emph{adına} sahiptir. Bu, $\Th{Q}$ dilindeki !!{formula}s{}in başka
!!{formula}s hakkında bir şeyler ``söyleyebileceği'' kavramsal bir çerçeve
sağlar. Aşağıdaki lemma sabit nokta lemması olarak bilinir.

\begin{lem}
$\Th{T}$, $\Th{Q}$'yu genişleten herhangi bir kuram ve $!B(x)$, yalnızca
$x$ değişkeni serbest olan herhangi bir !!{formula} olsun. O zaman
$\Th{T} \Proves !A \liff !B(\gn{!A})$ olacak biçimde !!a{sentence}~$!A$
vardır.
\end{lem}

Lemma, verilen herhangi bir $!B(x)$ özelliği için ``$!B(x)$ benim için
doğrudur'' diyen !!a{sentence}~$!A$ bulunduğunu ve $\Th{T}$'nin bunu
``bildiğini'' ileri sürer.

Böyle !!a{sentence} nasıl kurabiliriz? Epimenides paradoksunun Quine'a ait şu
biçimini ele alın:
\begin{quote}
``Önüne kendi tırnaklı biçimi getirildiğinde yanlışlık verir'' ifadesi,
önüne kendi tırnaklı biçimi getirildiğinde yanlışlık verir.
\end{quote}
Bu tümce doğrudan kendisine gönderme yapmaz. Yalnızca tırnak işaretleri
arasındaki sözdizimsel nesneler hakkında bir şey ileri sürer ve bu bakımdan
şu tümcelerle aynı düzeydedir:
\begin{enumerate}
\item ``Robert'' güzel bir addır.
\item ``Koştum.'' kısa bir tümcedir.
\item ``Üç sözcüğü vardır'' ifadesinin üç sözcüğü vardır.
\end{enumerate}
Peki ``önüne kendi tırnaklı biçimi getirildiğinde yanlışlık verir'' ifadesinin
önüne, kendisinin tırnak içine alınmış biçimini getirirsek ne olur? Özgün
tümceyi yeniden elde ederiz!{} Kısacası tümce, kendisinin yanlış olduğunu
ileri sürer.

